\documentclass[11pt,a4paper]{article}
\usepackage[margin=1in]{geometry}
\usepackage{natbib}
\usepackage{hyperref}
\usepackage{booktabs}
\usepackage{enumitem}
\usepackage{xcolor}
\usepackage{titlesec}

\bibliographystyle{plainnat}
\setcitestyle{authoryear,round}

\titleformat{\section}{\large\bfseries}{\thesection}{1em}{}
\titleformat{\subsection}{\normalsize\bfseries}{\thesubsection}{1em}{}

\title{\textbf{Research Gaps in World Models: A Survey of\\Open Problems and Opportunities}\\[0.5em]\large Based on a 61-Paper Literature Review Corpus}
\author{Rohit Lal}
\date{April 2026}

\begin{document}
\maketitle

\begin{abstract}
World models---learned internal representations of environment dynamics---are central to the next generation of autonomous agents. This report synthesizes findings from 61 papers spanning 2018--2026, covering latent dynamics for RL (Dreamer, TD-MPC), video generation as world simulation (GAIA-1, Cosmos, DIAMOND), Joint Embedding Predictive Architectures (I-JEPA, V-JEPA, LeWorldModel), and foundation world models. We identify ten major research gaps where significant contributions remain possible, supported by concrete evidence from the literature. For each gap, we articulate why the problem is important, what has been tried, and where opportunities for novel contributions exist.
\end{abstract}

\tableofcontents
\newpage

\section{Introduction}

A world model is a learned simulator of environment dynamics that allows an agent to predict the consequences of its actions without executing them in the real world \citep{ha2018worldmodels, lecun2022path}. The appeal is sample efficiency: instead of millions of real interactions, an agent trains inside imagined rollouts generated by the world model \citep{hafner2021dreamerv2, hafner2023dreamerv3}.

The field has matured rapidly. From Ha \& Schmidhuber's 2018 prototype to DreamerV3's domain-general agent \citep{hafner2023dreamerv3} to NVIDIA's billion-parameter Cosmos foundation model \citep{agarwal2025cosmos}, world models have evolved from toy demonstrations to systems that collect diamonds in Minecraft, drive cars in simulation, and control robots zero-shot. Yet fundamental gaps persist. This report identifies ten open problems where meaningful contributions are achievable, drawn from a systematic analysis of 61 papers.

\section{Gap 1: Long-Horizon Planning Under Error Accumulation}

\subsection{Why It Matters}
Compounding prediction errors over many steps is the single most cited limitation across the corpus. Nearly every world model paper acknowledges that prediction accuracy degrades with horizon length, limiting planning to short-horizon tasks.

\subsection{Current State}
\citet{maes2026leworldmodel} explicitly identifies short planning horizons as a key limitation of LeWorldModel. \citet{assran2025vjepa2} reports V-JEPA~2 is limited to $\sim$16-second predictions before error accumulation degrades accuracy. \citet{zhang2026hierarchical} demonstrates that flat (single-level) planning achieves 0\% success on non-greedy robotic pick-and-place tasks, while their hierarchical approach reaches 70\%---but still degrades at the longest horizons (61\% at $d{=}75$ on Push-T). \citet{hauri2026dreamercdp} evaluates only on Crafter, leaving long-horizon generalization untested.

\subsection{Opportunities}
\begin{itemize}[noitemsep]
    \item \textbf{Deeper hierarchies}: \citet{zhang2026hierarchical} explores only two-level hierarchy; LeCun's H-JEPA proposal \citep{lecun2022path} envisions multiple temporal scales, but no implementation exists beyond two levels.
    \item \textbf{Bidirectional planning feedback}: Current hierarchical methods are strictly top-down---if a high-level subgoal is unreachable, the system fails silently \citep{zhang2026hierarchical}. Bidirectional optimization between hierarchy levels is unexplored.
    \item \textbf{Adaptive horizon selection}: No existing method dynamically adjusts the planning horizon based on prediction confidence.
\end{itemize}

\section{Gap 2: End-to-End Training vs. Frozen Encoders}

\subsection{Why It Matters}
JEPA-based world models face a fundamental dilemma: freezing a pretrained encoder prevents representation collapse but limits adaptation to the dynamics task. End-to-end training allows richer representations but risks collapse.

\subsection{Current State}
\citet{maes2026leworldmodel} is the first JEPA world model to train stably end-to-end from pixels, using SIGReg \citep{balestriero2025lejepa} as the sole anti-collapse mechanism. However, \citet{terver2025drives} studies JEPA world models only in the frozen-encoder regime, explicitly noting that ``end-to-end training may alter which design choices matter most.'' \citet{nam2026causaljepa} and \citet{garrido2026learning} both freeze their encoders. \citet{assran2025vjepa2} freezes the V-JEPA~2 encoder during action-conditioned training, noting that representations are ``not optimized for control.''

\subsection{Opportunities}
\begin{itemize}[noitemsep]
    \item \textbf{Systematic comparison}: No paper has systematically compared frozen vs.\ end-to-end encoders across a diverse task suite with matched architectures.
    \item \textbf{Collapse-prevention landscape}: Beyond SIGReg \citep{balestriero2025lejepa}, KerJEPA \citep{zimmermann2025kerjepa} proposes kernel-based alternatives, and Barlow Twins is used in R2-Dreamer \citep{morihira2026r2dreamer}. A unified study comparing these anti-collapse mechanisms in the world model setting is missing.
    \item \textbf{Gradual unfreezing}: Progressive fine-tuning of frozen encoders during world model training---common in NLP---is unexplored for JEPA world models.
\end{itemize}

\section{Gap 3: Evaluation of World Model Quality}

\subsection{Why It Matters}
Standard metrics (next-token accuracy, task reward, FID/FVD) are structurally insensitive to world model quality. A model can achieve perfect accuracy while internalizing no valid world model.

\subsection{Current State}
\citet{vafa2024evaluating} demonstrates this sharply: a transformer with 1.00 next-token accuracy on NYC routing has sequence compression precision of 0.10; GPT-4 achieves perfect logic-puzzle accuracy with compression precision 0.21. Their Myhill-Nerode framework provides theoretically grounded behavioral metrics, but requires knowing the ground-truth world state space---limiting applicability to formal domains. \citet{ding2024understanding} documents that even Sora-class models fail basic physics tests.

\subsection{Opportunities}
\begin{itemize}[noitemsep]
    \item \textbf{Extending to continuous domains}: The Myhill-Nerode framework assumes a finite-state DFA. Extending it to continuous dynamics (physical simulations, robotic manipulation) is an open theoretical problem \citep{vafa2024evaluating}.
    \item \textbf{Characterizing the learning transition}: What training conditions cause a model to transition from superficial pattern matching to genuine world modeling? No paper addresses this.
    \item \textbf{Lightweight proxy metrics}: The compression/distinction metrics require expensive sampling. Developing efficient proxy metrics that correlate with world model quality would be practically valuable.
\end{itemize}

\section{Gap 4: Object-Centric Representation Scalability}

\subsection{Why It Matters}
Object-centric world models reason about individual entities, enabling compositional generalization and causal reasoning. But current methods struggle with real-world visual complexity.

\subsection{Current State}
\citet{nam2026causaljepa} proposes Causal-JEPA with object-level masking as latent interventions, improving counterfactual reasoning by $\sim$20\% on CLEVRER---but only on synthetic scenes. \citet{mosbach2025sold} outperforms DreamerV3 on multi-object manipulation (74.6\% vs. 33.1\%) but requires SAVi pretraining on 1M frames and uses a fixed slot count per environment. \citet{daniel2026lpwm} discovers keypoints and masks without supervision (ICLR 2026 Oral) but is limited to 128$\times$128 resolution with implicit tracking that breaks for long-range object traversals.

\subsection{Opportunities}
\begin{itemize}[noitemsep]
    \item \textbf{Real-world object-centric world models}: All current methods are validated on synthetic or simple tabletop scenes. Scaling to cluttered, real-world environments with occlusion, deformable objects, and varying object counts is the critical bottleneck.
    \item \textbf{Dynamic slot/particle allocation}: Both SOLD \citep{mosbach2025sold} and LPWM \citep{daniel2026lpwm} use fixed entity counts per environment. Adaptive allocation based on scene complexity is needed.
    \item \textbf{Causal validation}: \citet{nam2026causaljepa} proves that object masking induces causal inductive bias theoretically but does not validate on datasets with explicit causal graphs.
\end{itemize}

\section{Gap 5: Reconstruction-Free World Models}

\subsection{Why It Matters}
Pixel reconstruction wastes model capacity on task-irrelevant visual details and is computationally expensive. A growing body of work shows it may be unnecessary.

\subsection{Current State}
Three independent groups have converged on reconstruction-free world models: \citet{maes2026leworldmodel} (LeWorldModel, JEPA + SIGReg), \citet{morihira2026r2dreamer} (R2-Dreamer, Barlow Twins), and \citet{hauri2026dreamercdp} (Dreamer-CDP, JEPA-style cosine predictor). \citet{bredis2026nedreamer} (NE-Dreamer) shows that next-embedding prediction outperforms DreamerV3 on memory-demanding DMLab tasks. However, \citet{morihira2026r2dreamer} notes marginal gains on standard benchmarks (DMC, Meta-World), with clear advantages only on DMC-Subtle. \citet{hauri2026dreamercdp} evaluates on a single benchmark (Crafter).

\subsection{Opportunities}
\begin{itemize}[noitemsep]
    \item \textbf{Unified benchmark comparison}: No paper compares LeWorldModel, R2-Dreamer, Dreamer-CDP, and NE-Dreamer on the same tasks. A fair comparison across Atari, DMC, Crafter, and robotic manipulation is needed.
    \item \textbf{When does reconstruction help?}: Identifying the precise conditions under which pixel reconstruction provides useful gradients (vs. noise) would clarify when to use each paradigm.
    \item \textbf{Discrete action domains}: R2-Dreamer and Dreamer-CDP are only evaluated on continuous control. Whether reconstruction-free methods generalize to Atari-style discrete environments is untested \citep{morihira2026r2dreamer}.
\end{itemize}

\section{Gap 6: Representation-Action Alignment for Planning}

\subsection{Why It Matters}
World models learn to predict---but planning requires representations that are geometrically aligned with optimal control. A good predictor is not necessarily a good planner.

\subsection{Current State}
\citet{destrade2025valueguided} proposes shaping JEPA representations with IQL value functions so that Euclidean distance approximates goal-conditioned value. This improves planning success (0.71 vs. 0.55 on Wall-Small) but is only evaluated on simple 2D environments. \citet{kobanda2026intrinsic} proves theoretically that intrinsic-energy JEPAs induce quasimetric spaces aligned with optimal cost-to-go, but provides no experiments. \citet{parthasarathy2025closing} identifies the train-test gap: world models trained on prediction objectives perform poorly when used for gradient-based planning, and proposes online/adversarial fine-tuning.

\subsection{Opportunities}
\begin{itemize}[noitemsep]
    \item \textbf{Empirical validation of quasimetric JEPA}: \citet{kobanda2026intrinsic}'s theoretical connection between JEPA energies and quasimetric RL has zero experimental validation. Implementing and testing this is a direct contribution opportunity.
    \item \textbf{Joint representation-planning optimization}: Training world model representations to be simultaneously good for prediction and planning, rather than using separate training stages \citep{destrade2025valueguided}.
    \item \textbf{Complex-environment value shaping}: Scaling value-guided representation learning from 2D navigation to high-dimensional manipulation and 3D environments.
\end{itemize}

\section{Gap 7: Scaling Laws for World Models}

\subsection{Why It Matters}
LLMs benefit from well-characterized scaling laws (Chinchilla, etc.) that guide compute allocation. No equivalent exists for world models.

\subsection{Current State}
\citet{hansen2024tdmpc2} shows log-linear performance scaling from 1M to 317M parameters in continuous control. \citet{hafner2023dreamerv3} shows monotonic improvement from 12M to 400M parameters. \citet{assran2025vjepa2} scales to 1B but notes the ``scale ceiling is unknown.'' But systematic scaling laws relating compute, data, and world model capabilities (causal reasoning, counterfactual accuracy, planning horizon) have not been derived.

\subsection{Opportunities}
\begin{itemize}[noitemsep]
    \item \textbf{Chinchilla-style scaling analysis}: Controlled experiments varying model size, dataset size, and training compute to derive predictive scaling laws for world model quality metrics.
    \item \textbf{Capability-specific scaling}: Do different capabilities (next-state prediction, long-horizon accuracy, causal reasoning) scale at different rates? This fine-grained analysis is absent.
    \item \textbf{Compute-optimal allocation}: \citet{li2025rethinking} shows that in JEPA video SSL, compute should overwhelmingly favor the student over the teacher. Similar allocation principles for world model training are unknown.
\end{itemize}

\section{Gap 8: Physics, Causality, and Common Sense}

\subsection{Why It Matters}
Current world models learn correlational patterns from video data but do not reliably capture physical laws---gravity, conservation, fluid dynamics, object permanence.

\subsection{Current State}
\citet{ding2024understanding} documents systematic physics failures in Sora-class models. \citet{kong2025survey3d4d} argues that genuine world understanding requires 3D structure and 4D dynamics---capabilities current video generators approximate but do not achieve. \citet{nam2026causaljepa} introduces causal inductive bias via object masking but validates only on synthetic rigid-body physics.

\subsection{Opportunities}
\begin{itemize}[noitemsep]
    \item \textbf{Physics-informed world models}: Embedding PDE constraints, conservation laws, or Lagrangian mechanics into world model architectures. This is nascent and largely unexplored in the JEPA/Dreamer lineage.
    \item \textbf{Causal discovery from world models}: Using learned world models to extract causal graphs of environment dynamics, not just predict next states.
    \item \textbf{3D-native world models}: Most world models operate on 2D pixels. Building world models that natively reason in 3D (occupancy, point clouds, NeRF-like representations) is critical for robotics \citep{kong2025survey3d4d, zhu2026self}.
\end{itemize}

\section{Gap 9: Action Representation and Transfer}

\subsection{Why It Matters}
Most world models require explicit action labels for training, limiting their applicability to settings with instrumented agents. Learning from observation-only data (videos without actions) is essential for scaling.

\subsection{Current State}
\citet{garrido2026learning} shows that continuous constrained latent actions capture real-world complexity better than quantization, enabling world model training from unlabeled video. However, without a common embodiment, learned actions are camera-relative and spatially localized. Planning still requires a separately trained controller that maps known actions to latent space. \citet{maes2026leworldmodel} requires action labels; removing this dependency via inverse dynamics modeling is identified as future work.

\subsection{Opportunities}
\begin{itemize}[noitemsep]
    \item \textbf{Camera-invariant action representations}: Learning latent actions that transfer across viewpoints and embodiments, enabling cross-video and cross-robot transfer.
    \item \textbf{Direct latent-space planning}: Planning directly in the latent action space without a controller bridge remains an open problem due to challenges in sampling from the continuous latent distribution \citep{garrido2026learning}.
    \item \textbf{Language-conditioned actions}: \citet{assran2025vjepa2} notes that language-based goal specification would be ``far more practical'' than image goals. Bridging world models with language models for task specification is largely unaddressed.
\end{itemize}

\section{Gap 10: Stochastic and Multi-Modal Dynamics}

\subsection{Why It Matters}
Real environments are stochastic---the same action from the same state can lead to different outcomes. Most JEPA-based world models learn deterministic forward models, limiting their ability to capture uncertainty.

\subsection{Current State}
DreamerV3's RSSM explicitly models stochastic transitions via categorical latent variables \citep{hafner2023dreamerv3}. However, JEPA methods (LeWorldModel, TD-JEPA, Causal-JEPA) use deterministic predictors. \citet{mosbach2025sold} notes that SOLD's deterministic world model is a ``drawback in stochastic/unpredictable environments.'' \citet{gogl2026varjepa} and \citet{huang2026vjepa} propose variational JEPA formulations but are validated only on tabular data and simple domains respectively.

\subsection{Opportunities}
\begin{itemize}[noitemsep]
    \item \textbf{Stochastic JEPA world models}: Combining the prediction-in-representation-space principle of JEPA with stochastic latent variables to capture environment uncertainty. Var-JEPA \citep{gogl2026varjepa} and VJEPA \citep{huang2026vjepa} are first steps but need scaling to complex environments.
    \item \textbf{Multi-modal prediction}: World models that can represent branching futures (e.g., ``the ball might go left or right'') rather than averaging over modes. \citet{daniel2026lpwm} notes that unimodal Gaussian priors per particle may miss complex action distributions.
    \item \textbf{Risk-aware planning}: Using stochastic world models for robust planning under uncertainty, connecting to distributionally robust optimization.
\end{itemize}

\section{Summary of Opportunities}

Table~\ref{tab:gaps} summarizes the ten research gaps, their importance, and the most promising contribution directions.

\begin{table}[h]
\centering
\small
\caption{Summary of research gaps in world models with evidence and opportunity assessment.}
\label{tab:gaps}
\begin{tabular}{@{}p{3.5cm}p{4cm}p{4.5cm}@{}}
\toprule
\textbf{Research Gap} & \textbf{Key Evidence} & \textbf{Highest-Impact Opportunity} \\
\midrule
Long-horizon planning & 0\% flat vs. 70\% hierarchical \citep{zhang2026hierarchical} & Deeper hierarchies with bidirectional feedback \\
End-to-end vs. frozen & Only LeWM trains E2E \citep{maes2026leworldmodel} & Systematic frozen vs. E2E comparison \\
Evaluation metrics & 1.00 accuracy $\neq$ world model \citep{vafa2024evaluating} & Continuous-domain behavioral metrics \\
Object-centric scaling & Synthetic-only validation \citep{nam2026causaljepa} & Real-world object-centric world models \\
Reconstruction-free WM & 3 independent validations & Unified cross-method benchmark \\
Repr.-action alignment & Theory without experiments \citep{kobanda2026intrinsic} & Empirical quasimetric JEPA validation \\
Scaling laws & No Chinchilla-equivalent exists & Capability-specific scaling analysis \\
Physics \& causality & Sora fails physics \citep{ding2024understanding} & Physics-informed JEPA architectures \\
Action representation & Camera-relative only \citep{garrido2026learning} & Camera-invariant latent actions \\
Stochastic dynamics & JEPA models are deterministic & Stochastic JEPA world models \\
\bottomrule
\end{tabular}
\end{table}

\section{Conclusion}

The world model landscape in 2026 is rich with progress---from DreamerV3's domain generality to V-JEPA~2's zero-shot robotic control to LeWorldModel's elegant two-loss formulation---but the gaps identified here represent genuine open problems where novel contributions would advance the field. The most impactful directions are: (1) scaling hierarchical planning beyond two levels with bidirectional feedback, (2) systematically resolving the frozen-vs-end-to-end encoder question, (3) developing world model evaluation metrics for continuous domains, and (4) bridging the theory-practice gap in representation-action alignment for planning.

\bibliography{references}

\end{document}
